\documentclass[12pt]{article}
\usepackage{geometry}
\geometry{letterpaper, margin=1in}
\usepackage{hyperref}
\usepackage[usenames,dvipsnames]{xcolor}
\usepackage{microtype}
\usepackage{enumitem}

\definecolor{accent}{RGB}{26, 60, 110}

\begin{document}

%----------HEADER----------
\begin{center}
  {\LARGE\textbf{\textcolor{accent}{SITONG LI}}} \\[4pt]
  \small Ann Arbor, MI
  \;\textbullet\;
  \href{mailto:ruc.sitongli@gmail.com}{ruc.sitongli@gmail.com}
  \;\textbullet\;
  (952) 201-4026
\end{center}

\vspace{1em}

\noindent May 22, 2026

\vspace{1em}

\noindent Dear Prof.\ Pande, Prof.\ Schaner, Prof.\ Troyer-Moore, and the Yale Inclusion Economics Hiring Committee,

\vspace{1em}

I am writing to apply for the Research Associate / Predoctoral Fellow position supporting Yale Inclusion Economics' Kenya-based portfolio of applied development economics research. I am currently a full-time predoctoral research fellow at the Ross School of Business, University of Michigan, where I build and manage large-scale data infrastructure for applied economics research. The opportunity to bring these skills --- cleaning and managing complex administrative datasets, conducting rigorous statistical analyses, and supporting the full research production cycle --- to YIE's work evaluating digital empowerment interventions for women in East Africa is exactly where I want to direct my career, and I am eager to contribute.

At Michigan Ross, I work full-time as a predoctoral research fellow under faculty at Michigan, Yale, and Barnard. My core work is cleaning and managing large, heterogeneous administrative datasets: I built a scalable entity resolution pipeline integrating LLMs, fuzzy string matching, and word embeddings to link student application records across multiple institutional sources, and designed a computer vision pipeline for structured, large-scale demographic analysis. Working in a full-time, embedded research role --- with responsibility for the data infrastructure that faculty depend on to produce results --- has sharpened my ability to independently manage ambiguous data problems, maintain rigorous documentation, and deliver reliable outputs under academic timelines. It has also confirmed my commitment to the research path: I am actively preparing for a Ph.D.\ in economics, and a position supporting cutting-edge development economics research is the strongest foundation I can build toward that goal.

My most relevant methodological experience is from a research assistantship with Prof.\ Haiwen Zhang at the University of Minnesota on a project analyzing regulatory enforcement by the U.S.\ Environmental Protection Agency. I converted ten years of scanned regulatory correspondence into structured, analysis-ready datasets using OCR and regular expressions, managing over 100,000 documents across 1,000+ regulated entities. I then cleaned and merged those records with EPA administrative enforcement data and conducted probit regression analyses to examine the causal determinants of enforcement actions. This project --- converting messy, real-world administrative records into panel data suitable for causal inference, in service of a policy-relevant research question --- is closely analogous to the data-intensive survey and monitoring data work described in this position. In a concurrent project with Prof.\ Zhang, I validated empirical findings through panel regressions and event study models estimated in Stata across 10,000+ corporate filings, and as a research assistant to Prof.\ William Cassidy at Washington University in St.\ Louis, I conducted regression analyses linking macroeconomic outcomes to patterns documented in 70 years of congressional testimony. I am proficient in Stata, R, and Python, and am comfortable with the full workflow from raw data to econometric results.

I am also a strong written communicator. Across my research assistant positions, I have prepared analytical summaries, structured data documentation, and contributed to draft research materials for both academic and practitioner audiences. I am prepared to support literature reviews, background research, grant proposals, project reports, and research protocols of the kind this position requires. I am enthusiastic about the field research dimension of this role as well. I understand the position involves travel to study sites across Kenya and close collaboration with field research teams and implementing partner organizations, and I welcome the opportunity to work in that context. I am committed to the mission of policy-relevant development economics research conducted rigorously in the field, and I am prepared to contribute effectively in the environments that meaningful East Africa field work requires.

YIE's research agenda --- evaluating interventions to empower women as digital actors and ensuring digital programs are gender-intentional --- sits at the intersection of development economics and public policy where I believe the most important applied work is being done. I would be honored to contribute to it. I look forward to the opportunity to discuss how my background fits the needs of Prof.\ Pande's team and YIE's Kenya portfolio.

\vspace{1em}

\noindent Sincerely,\\[0.5em]
\textbf{Sitong Li}

\end{document}
